\documentclass[11pt]{article}

% Change "review" to "final" to generate the final (sometimes called camera-ready) version.
% Change to "preprint" to generate a non-anonymous version with page numbers.
\usepackage[review]{acl}

% Standard package includes
\usepackage{times}
\usepackage{latexsym}

% For proper rendering and hyphenation of words containing Latin characters (including in bib files)
\usepackage[T1]{fontenc}
% For Vietnamese characters
% \usepackage[T5]{fontenc}
% See https://www.latex-project.org/help/documentation/encguide.pdf for other character sets

% This assumes your files are encoded as UTF8
\usepackage[utf8]{inputenc}

% This is not strictly necessary, and may be commented out,
% but it will improve the layout of the manuscript,
% and will typically save some space.
\usepackage{microtype}

% This is also not strictly necessary, and may be commented out.
% However, it will improve the aesthetics of text in
% the typewriter font.
\usepackage{inconsolata}

%Including images in your LaTeX document requires adding
%additional package(s)
\usepackage{graphicx}

\usepackage{amsmath}
\usepackage{booktabs}
\usepackage{multirow}

% If the title and author information does not fit in the area allocated, uncomment the following
%
%\setlength\titlebox{<dim>}
%
% and set <dim> to something 5cm or larger.

\title{Decoupling Generation and Selection for Budget-Constrained and Faithful Summarization}

\author{
\textbf{Zeyu Wang} and
\textbf{Guanghua Wang}\thanks{Corresponding author.}
\\
\ Kean University, USA
\\
\small{
\texttt{\{wangzeyu, guanghua.wang\}@kean.edu}
}
}

\begin{document}
\maketitle
\begin{abstract}
This document is a supplement to the general instructions for *ACL authors. It contains instructions for using the \LaTeX{} style files for ACL conferences.
The document itself conforms to its own specifications, and is therefore an example of what your manuscript should look like.
These instructions should be used both for papers submitted for review and for final versions of accepted papers.
\end{abstract}
\section{Introduction}
\section{Related Work}
\subsection{Abstractive Summarization with Language Models}
\subsection{Faithfulness and Hallucination Reduction}
\subsection{Combinatorial Optimization for Summarization}
\section{Method}

\subsection{Problem Definition and Framework Overview}

Given an input document or document set $\mathcal{D}$, abstractive summarization aims to produce a concise and faithful textual summary that captures the salient information of $\mathcal{D}$. Most existing approaches formulate this task as conditional sequence generation, which achieves strong fluency but provides limited control over output length and factual consistency.

We instead cast abstractive summarization as a \emph{budget-constrained combinatorial selection problem} over a discrete candidate space. Specifically, we construct a candidate pool of sentences
$\mathcal{C} = \{c_1, \dots, c_n\}$
using a pretrained language model. The objective is to select a subset $S \subseteq \mathcal{C}$ that maximizes summary quality under an explicit budget constraint:
\begin{equation}
\max_{S \subseteq \mathcal{C}} f(S)
\quad \text{s.t.} \quad
\sum_{c_i \in S} l_i \le B,
\end{equation}
where $l_i$ denotes the cost (e.g., length) of candidate $c_i$, and $B$ is the budget.

We design $f(S)$ to capture the key properties of high-quality summaries:
\begin{itemize}
    \item \textbf{Utility (U):} promotes relevance and semantic coverage of the source document;
    \item \textbf{Factuality (F):} encourages consistency between candidate sentences and $\mathcal{D}$;
    \item \textbf{Redundancy (R):} penalizes semantic overlap among selected candidates to encourage diversity.
\end{itemize}
These components are integrated within a weighted combinatorial objective to guide subset selection under the budget constraint. In the current implementation, $f(S)$ consists of modular sentence utility and pairwise redundancy regularization. The utility term scores individual candidate sentences for coverage and factuality, while the redundancy term discourages selecting highly overlapping sentence pairs. This formulation matches the ILP, MMR, and DPP-style selectors used below without requiring a monotonicity assumption.

Figure~\ref{fig:pipeline} illustrates the overall framework, which consists of candidate generation, where a language model produces a diverse set of candidate sentences; combinatorial content planning, where an optimizer selects a subset of candidates by jointly considering utility, factuality, and redundancy under a budget constraint; and constrained summary realization, where the selected candidates are ordered to produce the final summary.

\begin{figure*}[t]
\centering
\includegraphics[width=\textwidth]{latex/Pipeline.png}
\caption{Overview of the proposed framework. A language model first generates a pool of candidate sentences. A combinatorial optimizer then selects a subset under a budget constraint by jointly considering utility, factuality, and redundancy. Finally, the selected content is ordered to produce the final summary.}
\label{fig:pipeline}
\end{figure*}

This formulation offers several advantages. It provides explicit control over the number of selected sentences in the current implementation; token- or word-level budgets can be incorporated by setting $l_i$ to the corresponding sentence length. It also improves \emph{faithfulness} by restricting the final output to content that is explicitly evaluated against the source document. Finally, it enables a \emph{modular integration} of neural generation and combinatorial optimization, allowing pretrained language models to be used without retraining while enforcing global constraints explicitly at the selection stage.




\subsection{Candidate Generation}

The first stage converts the continuous generation space of a pretrained summarizer into a finite set of sentence candidates. For each input $\mathcal{D}$, a generator $G$ produces $K$ candidate summaries,
$Y=\{y_1,\ldots,y_K\}$. These candidates provide multiple views of the same source document: they often share central facts while varying in wording, ordering, and secondary details. Collecting them before selection gives the optimizer access to both recurring content and complementary generated information.

Each candidate summary is sentence-segmented, and exact duplicate sentences are removed to form the candidate pool
$\mathcal{C}=\{c_1,\ldots,c_n\}$. We keep the provenance of each sentence, including the source beam index and within-summary sentence position, but selection is performed over the deduplicated sentence pool.

This design decouples surface realization from content selection. The language model proposes fluent abstractive sentence candidates, and the optimizer searches the resulting finite pool for a compact subset that best satisfies the budget, coverage, factuality, and redundancy criteria. The final summary is assembled from the selected candidate sentences after optimization. Since the selector only chooses from generated candidates, it cannot recover information that is absent from the candidate pool; generator quality and candidate diversity therefore define an important upper bound on the method.

\subsection{Scoring Functions}

For each candidate sentence $c_i$, we compute source coverage and factual consistency scores. For each pair $(c_i,c_j)$, we compute a redundancy score that estimates how much information would be repeated if both sentences were selected. The unified objective uses separate weights for coverage, factuality, and redundancy: coverage and factuality define sentence-level utility, while redundancy is applied as a pairwise selection penalty or diversity control term.

These scores make candidates from the same generated pool directly comparable. Sentences that are central to the source and well supported by the input should receive high utility, while fluent but weakly grounded or repetitive sentences should be down-weighted. Pairwise redundancy then controls how strongly the optimizer avoids selecting overlapping candidates. Thus, scoring provides both sentence-level evidence and subset-level diversity information for optimization.

\paragraph{Utility}
Coverage is measured with ROUGE recall against the source document:
\begin{equation}
g_i =
\operatorname{R}_{\mathrm{uni,recall}}(\mathcal{D}, c_i)
+
\operatorname{R}_{\mathrm{bi,recall}}(\mathcal{D}, c_i).
\end{equation}
This term favors sentences that reuse salient unigram and bigram content from the input.

\paragraph{Factuality}
Factuality is measured at the sentence level using MiniCheck. Given the source document and a candidate sentence, MiniCheck returns the probability that the sentence is consistent with the source:
\begin{equation}
m_i = P_{\mathrm{MiniCheck}}(\mathrm{consistent}\mid \mathcal{D}, c_i).
\end{equation}
The final sentence utility is a weighted sum of coverage and factuality:
\begin{equation}
u_i = \lambda_{\mathrm{cov}}g_i
    + \lambda_{\mathrm{fact}}m_i .
\end{equation}
This lets the selector prefer sentences that are both informative and supported by the source.

\paragraph{Redundancy}
Redundancy is represented as a pairwise similarity matrix $R\in\mathbb{R}^{n\times n}$. For $i\ne j$, we set
\begin{equation}
R_{ij} = \operatorname{ROUGE\mbox{-}L}_{F1}(c_i,c_j),
\end{equation}
and set $R_{ii}=1$. High $R_{ij}$ indicates that selecting both sentences is likely to repeat content. The redundancy weight $\lambda_{\mathrm{red}}$ controls how strongly this pairwise overlap affects subset selection.

\subsection{Budget-Constrained Optimization}

We instantiate the selection stage with ILP, MMR, and DPP optimizers. In the current experiments, each sentence has unit cost, so the budget controls the maximum number of selected sentences. The budget is set separately for each dataset according to the expected summary granularity.

Given the candidate pool, the optimizer searches for a compact subset under the budget constraint. The goal is to keep source-supported and informative content while removing redundant or low-utility candidates.

\paragraph{ILP}
The ILP selector uses binary variables $x_i\in\{0,1\}$ to indicate whether sentence $c_i$ is selected. In the unified objective, we introduce auxiliary binary variables $y_{ij}$ for selected sentence pairs and solve:
\begin{equation}
\max_x
\sum_i u_i x_i
-
\alpha \sum_{i<j} R_{ij}y_{ij}
\end{equation}
subject to
\begin{equation}
1 \le \sum_i x_i \le B,
\end{equation}
and standard linearization constraints
$y_{ij}\le x_i$, $y_{ij}\le x_j$, and
$y_{ij}\ge x_i+x_j-1$.
The coefficient $\alpha$ is derived from the redundancy weight and budget. This soft-penalty formulation allows the optimizer to select fewer than $B$ sentences when every additional sentence would add more redundancy than utility.

\paragraph{MMR}
MMR selects sentences greedily. At each step, it chooses the remaining candidate with the best marginal score:
\begin{equation}
\operatorname{MMR}(c_i)=
\lambda u_i -
(1-\lambda)\max_{c_j\in S}R_{ij}.
\end{equation}
The first term rewards quality, while the second term penalizes similarity to already selected sentences. In our weighted selection setting, $\lambda$ is mapped from the redundancy weight so that stronger redundancy settings put more pressure on diversity.

\paragraph{DPP}
The DPP selector uses a DPP-inspired quality-weighted similarity kernel:
\begin{equation}
L = QSQ,
\end{equation}
where $Q=\operatorname{diag}(q_1,\ldots,q_n)$ contains clipped sentence utilities and $S$ is the redundancy-weighted similarity matrix. Since ROUGE-based pairwise similarity is a heuristic and does not by itself guarantee a valid DPP kernel, we use this selector as a deterministic diversity-oriented approximation. Greedy MAP inference is applied to select a subset with a high log-determinant score, favoring high-quality sentences while avoiding near-duplicate candidates.

\subsection{Summary Realization}

After selection, the final summary is obtained by ordering the chosen sentences and concatenating them. The default realization strategy is \emph{source-similarity ordering}: each selected sentence is matched to the source sentence with the highest ROUGE-L F1 score, and selected sentences are sorted by the matched source position. This produces a stable document-level order without asking the language model to rewrite the selected content. This realization step preserves the optimized subset directly, but it may also inherit discourse gaps when selected sentences originate from different generated summaries. Because the final summary is a concatenation of selected candidates, the sentence budget is preserved exactly at realization time.

\section{Experimental Setup}
\subsection{Datasets}
We evaluate on news summarization benchmarks. CNN/DailyMail \cite{nallapati-etal-2016-abstractive} is a single-document news summarization dataset; we use its standard test split. Multi-News \cite{fabbri-etal-2019-multi} is a multi-document summarization dataset, where each example contains a cluster of related news articles and a human-written synthesis summary.

\subsection{Baselines}
We compare our optimization variants against both direct generation baselines and prior summarization systems. The direct local baselines include BART \cite{lewis-etal-2020-bart}, PRIMERA \cite{xiao-etal-2022-primera}, Qwen3.5-9B \cite{qwenteam2026qwen35omnitechnicalreport}, Llama-3.1-8B-Instruct \cite{llama3modelcard}, and Gemma-4-E4B-it \cite{gemma2024, gemma4_e4b_it}. We separately treat prior factuality- and reranking-oriented systems as reference baselines, including FactEdit \cite{balachandran-etal-2022-correcting}, SimCLS \cite{liu-liu-2021-simcls}, BRIO \cite{liu-etal-2022-brio}, EFactSum \cite{dixit-etal-2023-improving}, and SummaReranker \cite{ravaut-etal-2022-summareranker}. When outputs from prior systems are not re-evaluated under our pipeline, their numbers are interpreted as external reference points rather than strictly controlled comparisons.

Our proposed systems are denoted as \textsc{Generator}+MMR, \textsc{Generator}+ILP, and \textsc{Generator}+DPP. On CNN/DailyMail, we apply these selectors to BART and Llama candidate pools. On Multi-News, we apply them to PRIMERA and Llama candidate pools. Selector comparisons are made within the same generator-specific candidate pool whenever possible, which separates the downstream selection behavior from the candidate generator.

\subsection{Implementation Details}
All experiments are implemented with Hugging Face Transformers and PyTorch. In optimization runs, we use multiple generated candidates to populate the sentence pool for each input, and the selectors operate on the extracted pool. For BART, we use \texttt{facebook/bart-large-cnn} with the checkpoint's official summarization settings, including its length control, no-repeat trigram blocking, and early stopping configuration. The direct BART baseline uses the checkpoint's standard beam setting, while optimization runs use a larger beam to expose more candidate sentences. We therefore interpret direct-generation and optimization results as a search-and-selection comparison, and selector-level comparisons are made most directly among methods that share the same candidate pool. For Multi-News, we use \texttt{allenai/PRIMERA-multinews}; the direct baseline follows the model's default beam setting, and optimization runs use a larger candidate pool unless otherwise specified.

For instruction-tuned language models, we use concise summary-only prompts with dataset-specific sentence ranges. The local checkpoints are \texttt{Qwen/Qwen3.5-9B}, \texttt{meta-llama/Llama-3.1-8B-Instruct}, and \texttt{google/gemma-4-E4B-it}. Direct baselines use the model's sampling-style decoding configuration, whereas combinatorial optimization runs request multiple candidates and then perform sentence-level selection. For Multi-News inputs, articles are kept in their dataset order, separated by document delimiters, and truncated only when required by the model-specific input cap.

The optimization stage uses the weighted joint objective described in the Method section. Before combining metrics, coverage, factuality, and redundancy scores are normalized separately within each candidate pool. For a sentence-level score such as coverage or factuality, we compute the minimum and maximum over all candidate sentences in the pool, subtract the minimum from each score, and divide by the resulting score range. For redundancy, we apply the same min--max scaling to the off-diagonal entries of the pairwise redundancy matrix and keep the diagonal entries as self-similarity values. This places all components on a comparable scale while preserving their within-pool ordering. ILP is solved with the CBC backend through PuLP. DPP uses greedy MAP inference over the constructed kernel. MMR uses deterministic greedy selection. For fair comparison among ILP, DPP, and MMR, later selectors can reuse the same cached stage-one candidate pool produced by the first optimization run.

\subsection{Evaluation Metrics}
We report both content-overlap and factuality-oriented metrics. ROUGE is computed with the Hugging Face \texttt{evaluate} implementation using stemming and NLTK sentence splitting for ROUGE-Lsum. We treat summary-level ROUGE and bigram overlap as the main content-overlap indicators, and report unigram and longest-common-subsequence variants for completeness. We also report BERTScore F1 with \texttt{roberta-large}.

For faithfulness, we report FactCC, MiniCheck, AlignScore, and FactKB. These metrics are computed against the source document rather than the reference summary. Because MiniCheck is also used as a selection signal, MiniCheck evaluation is interpreted as an optimization-aligned diagnostic rather than a fully independent factuality measure. We therefore use FactCC, AlignScore, and FactKB as complementary external checks. Scores are reported as percentages, and higher values indicate better performance.
\section{Results}
\subsection{Main Results}
\begin{table*}[htbp]
\centering
\small
\renewcommand{\arraystretch}{1.1}
\setlength{\tabcolsep}{4.5pt}
\resizebox{\linewidth}{!}{%
\begin{tabular}{@{} l l cccc c cccc @{}}
\toprule
\textbf{Dataset} & \textbf{Method}
 & \multicolumn{4}{c}{\textbf{ROUGE} ($\uparrow$)}
 & \textbf{BERTScore} ($\uparrow$)
 & \multicolumn{4}{c}{\textbf{Faithfulness} ($\uparrow$)} \\
\cmidrule(lr){3-6} \cmidrule(lr){7-7} \cmidrule(l){8-11}
 &
 & R-1 & R-2 & R-L & R-Lsum
 & F1
 & FactCC & MiniCheck & AlignScore & FactKB \\
\midrule

\multirow{16}{*}{CNN/DM}
 & BART        & 44.05 & 21.07 & 30.62 & 40.99 & 88.21 & 75.83 & 94.90 & 91.56 & 98.49 \\
 & Qwen3.5-9B        & 34.88 & 10.91 & 20.33 & 30.77 & 86.96 & 40.16 & 56.82 & 65.17 & 98.12 \\
 & FactEdit        & 42.53 & 20.48 & 39.74 & 40.76 & 88.17 & 76.03 & 94.82 & 91.46 & 98.68 \\
 & SimCLS          & 46.67 & 22.15 & 43.54 & 43.09 & 66.14 & 59.49 & 83.34 & 82.37 & 80.75 \\
 & BRIO-Mul        & 47.78 & 23.55 & 44.57 & 44.12 & 87.84 & 54.06 & 84.67 & 74.48 & 85.20 \\
 & BRIO-Ctr        & 47.28 & 22.93 & 44.15 & 44.10 & 89.11 & 62.62 & 88.42 & 84.98 & 78.89 \\
 & EFactSum        & 44.37 & 21.28 & 40.92 & 41.07 & 88.36 & 60.74 & 88.33 & 82.02 & 91.38 \\
 & SummaReranker   & 46.62 & 22.39 & 43.59 & 43.80 & 88.47 & 68.17 & 93.81 & 88.27 & 98.10 \\
 & Llama-3-8B      & 38.98 & 14.82 & 24.20 & 35.28 & 87.72 & 46.78 & 76.00 & 75.08 & 98.51 \\
 & Gemma-4-E4B-it   & 30.51 & 6.41 & 18.07 & 27.40 & 86.19 & 53.99 & 69.49 & 71.32 & 89.59 \\
 \cmidrule(l){2-11}
 & BART+MMR            & 41.85 & 19.44 & 28.03 & 38.99 & 87.75 & 76.55 & 96.22 & 92.42 & 98.67 \\
 & BART+ILP             & 42.84 & 19.91 & 28.46 & 39.87 & 87.91 & 76.71 & 96.76 & 92.56 & 98.65 \\
 & BART+DPP             & 41.67 & 19.36 & 27.96 & 38.85 & 87.72 & 76.92 & 96.93 & 93.14 & 98.83 \\
 & Llama-3-8B+MMR            &  &  &  &  &  &  &  &  &  \\
 & Llama-3-8B+ILP            &  &  &  &  &  &  &  &  &  \\
 & Llama-3-8B+DPP            &  &  &  &  &  &  &  &  &  \\

\midrule

\multirow{16}{*}{Multi-News}
 & PRIMERA            &  &  &  &  &  &  &  &  &  \\
 & Qwen3.5-9B        &  &  &  &  &  &  &  &  &  \\
 & FactEdit        &  &  &  &  &  &  &  &  &  \\
 & SimCLS          &  &  &  &  &  &  &  &  &  \\
 & BRIO-Mul        &  &  &  &  &  &  &  &  &  \\
 & BRIO-Ctr        &  &  &  &  &  &  &  &  &  \\
 & EFactSum        &  &  &  &  &  &  &  &  &  \\
 & SummaReranker   &  &  &  &  &  &  &  &  &  \\
 & Llama-3-8B      &  &  &  &  &  &  &  &  &  \\
 & Gemma-4-E4B-it   &  &  &  &  &  &  &  &  &  \\
 \cmidrule(l){2-11}
 & PRIMERA+MMR        &  &  &  &  &  &  &  &  &  \\
 & PRIMERA+ILP        &  &  &  &  &  &  &  &  &  \\
 & PRIMERA+DPP        &  &  &  &  &  &  &  &  &  \\
 &Llama-3-8B+MMR          &  &  &  &  &  &  &  &  &  \\
 & Llama-3-8B+ILP          &  &  &  &  &  &  &  &  &  \\
 & Llama-3-8B+DPP          &  &  &  &  &  &  &  &  &  \\

\bottomrule
\end{tabular}%
}
\caption{Performance comparison on CNN/DailyMail and Multi-News datasets.}
\label{tab:main_results}
\end{table*}
\subsection{Ablation Study}

\section{Analysis}
\subsection{Trade-offs between Coverage, Redundancy, and Faithfulness}
\subsection{Effect of Candidate Pool Size}
\subsection{Case Studies}

\section{Limitations}
\section{Conclusion}




\section{Introduction}

These instructions are for authors submitting papers to *ACL conferences using \LaTeX. They are not self-contained. All authors must follow the general instructions for *ACL proceedings,\footnote{\url{http://acl-org.github.io/ACLPUB/formatting.html}} and this document contains additional instructions for the \LaTeX{} style files.

The templates include the \LaTeX{} source of this document (\texttt{acl\_latex.tex}),
the \LaTeX{} style file used to format it (\texttt{acl.sty}),
an ACL bibliography style (\texttt{acl\_natbib.bst}),
an example bibliography (\texttt{custom.bib}),
and the bibliography for the ACL Anthology (\texttt{anthology.bib}).

\section{Engines}

To produce a PDF file, pdf\LaTeX{} is strongly recommended (over original \LaTeX{} plus dvips+ps2pdf or dvipdf).
The style file \texttt{acl.sty} can also be used with
lua\LaTeX{} and
Xe\LaTeX{}, which are especially suitable for text in non-Latin scripts.
The file \texttt{acl\_lualatex.tex} in this repository provides
an example of how to use \texttt{acl.sty} with either
lua\LaTeX{} or
Xe\LaTeX{}.

\section{Preamble}

The first line of the file must be
\begin{quote}
\begin{verbatim}
\documentclass[11pt]{article}
\end{verbatim}
\end{quote}

To load the style file in the review version:
\begin{quote}
\begin{verbatim}
\usepackage[review]{acl}
\end{verbatim}
\end{quote}
For the final version, omit the \verb|review| option:
\begin{quote}
\begin{verbatim}
\usepackage{acl}
\end{verbatim}
\end{quote}

To use Times Roman, put the following in the preamble:
\begin{quote}
\begin{verbatim}
\usepackage{times}
\end{verbatim}
\end{quote}
(Alternatives like txfonts or newtx are also acceptable.)

Please see the \LaTeX{} source of this document for comments on other packages that may be useful.

Set the title and author using \verb|\title| and \verb|\author|. Within the author list, format multiple authors using \verb|\and| and \verb|\And| and \verb|\AND|; please see the \LaTeX{} source for examples.

By default, the box containing the title and author names is set to the minimum of 5 cm. If you need more space, include the following in the preamble:
\begin{quote}
\begin{verbatim}
\setlength\titlebox{<dim>}
\end{verbatim}
\end{quote}
where \verb|<dim>| is replaced with a length. Do not set this length smaller than 5 cm.

\section{Document Body}

\subsection{Footnotes}

Footnotes are inserted with the \verb|\footnote| command.\footnote{This is a footnote.}

\subsection{Tables and figures}

See Table~\ref{tab:accents} for an example of a table and its caption.
\textbf{Do not override the default caption sizes.}

\begin{table}
  \centering
  \begin{tabular}{lc}
    \hline
    \textbf{Command} & \textbf{Output} \\
    \hline
    \verb|{\"a}|     & {\"a}           \\
    \verb|{\^e}|     & {\^e}           \\
    \verb|{\`i}|     & {\`i}           \\
    \verb|{\.I}|     & {\.I}           \\
    \verb|{\o}|      & {\o}            \\
    \verb|{\'u}|     & {\'u}           \\
    \verb|{\aa}|     & {\aa}           \\\hline
  \end{tabular}
  \begin{tabular}{lc}
    \hline
    \textbf{Command} & \textbf{Output} \\
    \hline
    \verb|{\c c}|    & {\c c}          \\
    \verb|{\u g}|    & {\u g}          \\
    \verb|{\l}|      & {\l}            \\
    \verb|{\~n}|     & {\~n}           \\
    \verb|{\H o}|    & {\H o}          \\
    \verb|{\v r}|    & {\v r}          \\
    \verb|{\ss}|     & {\ss}           \\
    \hline
  \end{tabular}
  \caption{Example commands for accented characters, to be used in, \emph{e.g.}, Bib\TeX{} entries.}
  \label{tab:accents}
\end{table}

As much as possible, fonts in figures should conform
to the document fonts. See Figure~\ref{fig:experiments} for an example of a figure and its caption.

Using the \verb|graphicx| package graphics files can be included within figure
environment at an appropriate point within the text.
The \verb|graphicx| package supports various optional arguments to control the
appearance of the figure.
You must include it explicitly in the \LaTeX{} preamble (after the
\verb|\documentclass| declaration and before \verb|\begin{document}|) using
\verb|\usepackage{graphicx}|.

\begin{figure}[t]
  \includegraphics[width=\columnwidth]{example-image-golden}
  \caption{A figure with a caption that runs for more than one line.
    Example image is usually available through the \texttt{mwe} package
    without even mentioning it in the preamble.}
  \label{fig:experiments}
\end{figure}

\begin{figure*}[t]
  \includegraphics[width=0.48\linewidth]{example-image-a} \hfill
  \includegraphics[width=0.48\linewidth]{example-image-b}
  \caption {A minimal working example to demonstrate how to place
    two images side-by-side.}
\end{figure*}

\subsection{Hyperlinks}

Users of older versions of \LaTeX{} may encounter the following error during compilation:
\begin{quote}
\verb|\pdfendlink| ended up in different nesting level than \verb|\pdfstartlink|.
\end{quote}
This happens when pdf\LaTeX{} is used and a citation splits across a page boundary. The best way to fix this is to upgrade \LaTeX{} to 2018-12-01 or later.

\subsection{Citations}

\begin{table*}
  \centering
  \begin{tabular}{lll}
    \hline
    \textbf{Output}           & \textbf{natbib command} & \textbf{ACL only command} \\
    \hline
    \citep{Gusfield:97}       & \verb|\citep|           &                           \\
    \citealp{Gusfield:97}     & \verb|\citealp|         &                           \\
    \citet{Gusfield:97}       & \verb|\citet|           &                           \\
    \citeyearpar{Gusfield:97} & \verb|\citeyearpar|     &                           \\
    \citeposs{Gusfield:97}    &                         & \verb|\citeposs|          \\
    \hline
  \end{tabular}
  \caption{\label{citation-guide}
    Citation commands supported by the style file.
    The style is based on the natbib package and supports all natbib citation commands.
    It also supports commands defined in previous ACL style files for compatibility.
  }
\end{table*}

Table~\ref{citation-guide} shows the syntax supported by the style files.
We encourage you to use the natbib styles.
You can use the command \verb|\citet| (cite in text) to get ``author (year)'' citations, like this citation to a paper by \citet{Gusfield:97}.
You can use the command \verb|\citep| (cite in parentheses) to get ``(author, year)'' citations \citep{Gusfield:97}.
You can use the command \verb|\citealp| (alternative cite without parentheses) to get ``author, year'' citations, which is useful for using citations within parentheses (e.g. \citealp{Gusfield:97}).

A possessive citation can be made with the command \verb|\citeposs|.
This is not a standard natbib command, so it is generally not compatible
with other style files.

\subsection{References}

\nocite{Ando2005,andrew2007scalable,rasooli-tetrault-2015}

The \LaTeX{} and Bib\TeX{} style files provided roughly follow the American Psychological Association format.
If your own bib file is named \texttt{custom.bib}, then placing the following before any appendices in your \LaTeX{} file will generate the references section for you:
\begin{quote}
\begin{verbatim}
\bibliography{custom}
\end{verbatim}
\end{quote}

You can obtain the complete ACL Anthology as a Bib\TeX{} file from \url{https://aclweb.org/anthology/anthology.bib.gz}.
To include both the Anthology and your own .bib file, use the following instead of the above.
\begin{quote}
\begin{verbatim}
\bibliography{anthology,custom}
\end{verbatim}
\end{quote}

Please see Section~\ref{sec:bibtex} for information on preparing Bib\TeX{} files.

\subsection{Equations}

An example equation is shown below:
\begin{equation}
  \label{eq:example}
  A = \pi r^2
\end{equation}

Labels for equation numbers, sections, subsections, figures and tables
are all defined with the \verb|\label{label}| command and cross references
to them are made with the \verb|\ref{label}| command.

This an example cross-reference to Equation~\ref{eq:example}.

\subsection{Appendices}

Use \verb|\appendix| before any appendix section to switch the section numbering over to letters. See Appendix~\ref{sec:appendix} for an example.

\section{Bib\TeX{} Files}
\label{sec:bibtex}

Unicode cannot be used in Bib\TeX{} entries, and some ways of typing special characters can disrupt Bib\TeX's alphabetization. The recommended way of typing special characters is shown in Table~\ref{tab:accents}.

Please ensure that Bib\TeX{} records contain DOIs or URLs when possible, and for all the ACL materials that you reference.
Use the \verb|doi| field for DOIs and the \verb|url| field for URLs.
If a Bib\TeX{} entry has a URL or DOI field, the paper title in the references section will appear as a hyperlink to the paper, using the hyperref \LaTeX{} package.

\section*{Limitations}

This document does not cover the content requirements for ACL or any
other specific venue.  Check the author instructions for
information on
maximum page lengths, the required ``Limitations'' section,
and so on.

\section*{Acknowledgments}

This document has been adapted
by Steven Bethard, Ryan Cotterell and Rui Yan
from the instructions for earlier ACL and NAACL proceedings, including those for
ACL 2019 by Douwe Kiela and Ivan Vuli\'{c},
NAACL 2019 by Stephanie Lukin and Alla Roskovskaya,
ACL 2018 by Shay Cohen, Kevin Gimpel, and Wei Lu,
NAACL 2018 by Margaret Mitchell and Stephanie Lukin,
Bib\TeX{} suggestions for (NA)ACL 2017/2018 from Jason Eisner,
ACL 2017 by Dan Gildea and Min-Yen Kan,
NAACL 2017 by Margaret Mitchell,
ACL 2012 by Maggie Li and Michael White,
ACL 2010 by Jing-Shin Chang and Philipp Koehn,
ACL 2008 by Johanna D. Moore, Simone Teufel, James Allan, and Sadaoki Furui,
ACL 2005 by Hwee Tou Ng and Kemal Oflazer,
ACL 2002 by Eugene Charniak and Dekang Lin,
and earlier ACL and EACL formats written by several people, including
John Chen, Henry S. Thompson and Donald Walker.
Additional elements were taken from the formatting instructions of the \emph{International Joint Conference on Artificial Intelligence} and the \emph{Conference on Computer Vision and Pattern Recognition}.

% Bibliography entries for the entire Anthology, followed by custom entries
%\bibliography{anthology,custom}
% Custom bibliography entries only
\bibliography{custom}

\appendix

\section{Example Appendix}
\label{sec:appendix}

This is an appendix.

\end{document}
